% ============================================================================
% Fichier  : front/02_mode_emploi.tex
% Titre    : Mode d'emploi du cours
% Auteur   : MINKA MI NGUIDJOÏ Thierry Emmanuel
% Version  : 1.0 -- novembre 2026
% ============================================================================

\chapter*{Mode d'emploi du cours}
\addcontentsline{toc}{chapter}{Mode d'emploi du cours}
\markboth{Mode d'emploi}{Mode d'emploi}

\section*{Structure générale}

Le cours compte un prologue et six parties, soit 22~chapitres, pour
30~heures en présentiel réparties en quinze séances de deux heures, qui
couvrent l'ensemble des parties. Les ateliers de questions-réponses par
chapitre et les dix laboratoires de 60~minutes de NDAP SUKULU sont des
travaux à faire à la maison, hors de ces 30~heures. Les Parties~V et~VI
servent en outre de point de départ aux mémoires.

\begin{tableaucours}{L{1.6cm}L{4.4cm}L{4.4cm}c}{%
    \tableauHeader{Partie} &
    \tableauHeader{Thème} &
    \tableauHeader{Compétences visées} &
    \tableauHeader{Niveau}
}
Prologue & Quarante-cinq questions & Distinguer constatation, impression et
opinion & \niveaubase \\
I & Fondements de l'audit et du contrôle interne & Situer l'audit, évaluer
un contrôle, coter un risque, repérer la fraude & \niveaubase \\
II & Méthode de la mission & Planifier, décomposer, collecter la preuve,
conclure avec $\mu$, $U$ et $H$ & \niveaumoyen \\
III & Cadre normatif, juridique et déontologique & Choisir ses normes et
ses critères~; connaître ses obligations & \niveaubase \\
IV & Pratique de l'audit des SI & Auditer la gouvernance, une application,
un projet, la sécurité, la continuité & \niveaumoyen\,--\,\niveauexpert \\
V & Défis avancés & Audit continu, analyse de données, IA &
\niveauexpert \\
VI & Recherche~: l'auditologie & Lire et prolonger le programme de
recherche & \niveauexpert \\
\end{tableaucours}

\section*{Les encadrés}

\begin{boxdef}
\textbf{Définition}. Introduit un terme avec sa définition~; les termes
définis sont indexés et repris dans le glossaire (Annexe~\ref{ann:glossaire}).
\end{boxdef}

\begin{boxscenario}{Titre du scénario}
\textbf{Scénario opérationnel}. Un épisode d'une des trois opérations du
cours~: LIZBIZ, SANAGA, MEFOU.
\end{boxscenario}

\begin{boxlecture}{Titre de la lecture}
\textbf{Lecture auditologique}. Une cotation \OUI{}~/ \NON{}~/ \NEK{}, sa
constatation, son analyse avec un niveau de confiance, et les trois mesures
du domaine, rappelées en bas à droite.

\medskip
\hfill\lecture{0{,}65}{0{,}25}{1{,}24}
\end{boxlecture}

\begin{boxobjet}{Code · Intitulé}
Présente un contrôle sous une forme constante~: objectif, critère, question,
risque, conséquence, puis test, preuve et trace~; ou une observation en
sept éléments.
\end{boxobjet}

\begin{boxprocedure}
Dit ce que l'auditeur doit faire, et ce que le dossier doit contenir, pour
qu'une étape soit opposable.
\end{boxprocedure}

\begin{boxmetier}{Auditeur SI}
\textbf{Application métier}. Traduit le chapitre en gestes concrets pour
l'un des métiers visés~: auditeur SI, commissaire aux comptes, contrôleur
interne, responsable GRC, RSSI.
\end{boxmetier}

\begin{boxjuris}
Cite un texte de loi ou de règlement, avec son degré de vérification, puis
en tire la portée pour l'audit.
\end{boxjuris}

\begin{boxalert}
\textbf{Violation typique} et \textbf{règle}~: une erreur fréquente, et ce
qu'il faut faire à la place.
\end{boxalert}

\begin{boxresume}
\textbf{À retenir}. Les points essentiels du chapitre.
\end{boxresume}

\begin{boxexo}
\textbf{Exercice}. Gradué en trois niveaux~: \niveaubase\ fondamental,
\niveaumoyen\ intermédiaire, \niveauexpert\ expert.
\end{boxexo}

\section*{Conventions}

\begin{itemize}
    \item Les réponses s'écrivent \OUI{}, \NON{} et \NEK{}~; un contrôle sans
          objet est écarté, raison écrite, et n'entre pas dans les calculs.
    \item $\mu$ est le taux de conformité, $U$ le taux d'incertitude, $H$
          l'entropie en bits (au plus $\log_2 3 \approx 1{,}585$)~; la règle
          de décision examine $U$ d'abord, $\mu$ ensuite, $H$ pour contrôler.
    \item Les seuils du quadrant ($\mu = 0{,}5$~; $U = 0{,}20$) et la
          conversion en niveaux de maturité sont des conventions
          pédagogiques, déclarées comme telles.
    \item Les axiomes de l'auditologie portent la numérotation de l'article
          A0 (Chapitre~\ref{chap:auditologie}).
    \item Chaque analyse indique un niveau de confiance~: faible, moyenne ou
          élevée.
    \item Les \termecle{termes clés} apparaissent en gras~; les normes
          s'écrivent \norme{ISO/IEC~27001}, les lois camerounaises
          \loicam{n°~2010/012}.
\end{itemize}

\section*{NDAP SUKULU}

Le cours est accompagné, sous le code \codeUE, d'une fiche
d'auto-évaluation par chapitre (prologue compris) et de dix laboratoires
pratiques. Leur présentation et le tableau des labs figurent dans la section
« NDAP SUKULU, la salle de classe compagnon », qui suit ce mode d'emploi.

\section*{Matériel d'accompagnement}

\begin{itemize}
    \item \textbf{Annexes}~: glossaire, référentiel des outils, modèles et
          grilles, textes juridiques et normatifs, bibliographie commentée,
          contacts (Annexes~\ref{ann:glossaire} à~\ref{ann:contacts}).
    \item \textbf{Corrigés des labs}~: réservés au formateur, présentés en
          séance.
\end{itemize}
